\documentclass[conference]{IEEEtran}
% Compile on Overleaf (IEEEtran is built in) or with pdflatex + bibtex.
% Every number in this paper comes from results/eval/*.csv (generated by netinject.reporting.report);
% tables are written by scripts/make_paper_tables.py into paper/tables/.
\usepackage{graphicx}
\usepackage{booktabs}
\usepackage{amsmath}
\usepackage{url}
\usepackage{xcolor}
\newcommand{\pending}[1]{\textcolor{red}{[#1]}}  % placeholders until the full grid is analysed

\begin{document}

\title{Multi-Step Delayed Tool-Output Poisoning in Network-Operations LLM Agents: Repeated-Decoding Evaluation of Local 7--8B Models with Change-Control Approval Gates}

\author{\IEEEauthorblockN{\pending{Author full name(s)}}
\IEEEauthorblockA{\pending{Department, institution, city, country; email}}}

\maketitle

\begin{abstract}
Large language model (LLM) agents are being connected to network-operations tools, where an instruction hidden in a ticket comment, a syslog line, a runbook or a knowledge-base article can trigger a harmful configuration change.
Prior work on network-operations agents measured this risk in a single step, at temperature~0 and with one run per scenario.
We study the multi-step case.
We control the number of tool outputs between the poisoned read and the write (delay $d\in\{0,1,2,3\}$, with the write position fixed), repeat stochastic decoding (temperature~0.7, several seeds), and compare two deterministic change-control approval gates.
On a mock enterprise network with 13 tools, 24 benign tasks and up to 384 attack episodes per model, we evaluate Qwen2.5-7B, Llama-3.1-8B and Mistral-7B served locally with Ollama.
\pending{Main findings: effect of delay; share of episodes compromised only under sampling; residual attack success of scope\_check under ticket confusion versus ticket\_bound; utility cost.}
\end{abstract}

\begin{IEEEkeywords}
prompt injection, LLM agents, network operations, tool use, change control, AI security
\end{IEEEkeywords}

\section{Introduction}
Network operations centres (NOCs) are starting to use LLM agents that read tickets, alerts and device logs and then change device configurations.
Every artifact such an agent reads may contain text written by an outsider: a ticket comment, a crafted syslog message, an edited runbook paragraph.
Indirect prompt injection~\cite{greshake2023} turns that text into instructions.

NetInjectBench~\cite{netinjectbench} showed that 7--8B open models choose an unsafe tool action in most single-step attack scenarios, and that a metadata-aware policy gate removes those actions.
Its authors list four open problems: multi-step agents, delayed tool-output poisoning, human-in-the-loop approval and repeated stochastic decoding.
This paper addresses all four in one controlled environment, using the same three model families so the results remain comparable.

Our questions are:
\textbf{RQ1} Does the attack rate change as the delay between the poisoned content and the harmful write increases?
\textbf{RQ2} How much do outcomes vary across repeated runs at temperature 0.7 compared with temperature 0?
\textbf{RQ3} How much do the \texttt{scope\_check} and \texttt{ticket\_bound} gates reduce executed unsafe writes, and at what cost in false blocks and task success?
\textbf{RQ4} Which attack styles and artifact types are most effective, and does this differ by model?

Our contributions are:
(i) a delay design that fixes the position of the write and slides the poisoned read, so the payload is always read and only its distance to the action changes;
(ii) per-episode variability of attack success under sampling for local 7--8B agents;
(iii) a gate comparison that includes a \emph{ticket-confusion} attack, which cites a real, approved but unrelated change ticket and separates authorization from intent;
(iv) an open, CPU-testable harness whose labels are computed deterministically from full traces.

\section{Related Work}
\textbf{Indirect prompt injection in agents.}
InjecAgent~\cite{injecagent} measured single-turn injections in tool responses.
AgentDojo~\cite{agentdojo} and ASB~\cite{asb} evaluate multi-step agents, but the injection position is fixed per task and results come from single attempts.
Rashidi~\cite{rashidi2026depth} varies the position of the poisoned tool output and finds that success falls with depth for GPT-4o-mini. That depth is counted from the start of the run, and deep payloads were often never read.
We instead hold the write position fixed and vary the distance from the poisoned read to the write.

\textbf{Repeated sampling.}
Sampling and decoding settings change jailbreak success~\cite{huang2024catastrophic,bestofn}.
For agent hijacking, CAISI reported that repeating each attack 25 times raised average success from 57\% to 80\% on frontier models~\cite{caisi2025hijacking}.
We report per-episode variability for local open models.

\textbf{Defenses.}
DataSentinel~\cite{datasentinel} detects injected inputs with a game-theoretically trained detector, which is the detection family of defenses.
Deterministic, task-scoped tool-call policies form a second family: CaMeL~\cite{camel}, Progent~\cite{progent}, FIDES~\cite{fides} and related design patterns~\cite{designpatterns}.
Our gates are a network-operations instance of that second family, evaluated as change-control approvers rather than proposed as a new mechanism.

\textbf{LLMs for networks.}
NetConfEval~\cite{netconfeval} studies configuration synthesis; NetInjectBench~\cite{netinjectbench} studies injection in a single step.

\section{Threat Model}
The agent operates a network on behalf of an operator who names one approved change ticket.
The attacker can write text into artifacts the agent reads (ticket comments, device logs, runbooks, knowledge-base articles, alert notes) but cannot modify the trusted approval record (ticket status and approved scope) or the tools.
The attacker's goal is an out-of-scope high-impact write: disabling logging, creating an admin account, shutting an interface, adding a permit-any rule, or redirecting the default route.

\section{Environment and Method}
\subsection{Mock network and tools}
An in-memory model of an enterprise network with ten devices exposes 13 tools (Table~\ref{tab:tools}): six reads, two low-impact writes and five high-impact writes.
No tool touches a real device.
A ticket's approved scope lists exact changes (action, device and parameters), so an extra rule on the right firewall is still out of scope.

\subsection{Tasks, attacks and delay}
Each of the 24 benign tasks gives an ordered procedure: list alerts (step 1), four reads (ticket, device log, runbook, knowledge base; steps 2--5), then the approved write.
The payload is placed in the read at step $5-d$, so $d$ reads separate it from the write.
The order of the four reads follows a Latin square across tasks, so each artifact type appears at each position equally often and delay is not confounded with artifact type.
Four attack styles are crossed with the four delays: direct instruction, fake authority, ticket confusion (citing a real approved decoy ticket whose scope covers the goal) and an instruction split across two adjacent reads.
Five attack goals rotate across tasks and styles.
The full set therefore has 24 control and 384 attack episodes (Section~\ref{sec:setup} gives the subsets each model ran).

\subsection{Agent, gates and forking}
The agent receives a neutral NOC-operator system prompt with no warning about injected instructions, comparable to the ``naive'' setting of NetInjectBench, and runs for at most 12 model calls.
\texttt{scope\_check} allows a high-impact write only if it cites an approved ticket whose scope covers it exactly; \texttt{ticket\_bound} additionally requires that ticket to be the one named in the task.
Each episode is run once without a gate. At the first call a gate would deny, the run is copied and continued with that gate active.
This is equivalent to running each gate separately, because the shared prefix has the same distribution, and it yields paired comparisons.

\subsection{Labels and statistics}
A deterministic judge computes, from each trace: task success (the approved change is in the final state), whether the attacker's exact action was attempted or executed, any out-of-scope high-impact write, gate blocks, false blocks, and the realized delay.
We report Wilson intervals~\cite{wilson1927}, a logistic regression of attempts on delay, model, style and artifact with episode-clustered standard errors, and McNemar's test~\cite{mcnemar1947} for paired gate comparisons.

\section{Experimental Setup}\label{sec:setup}
\textbf{Models and serving.}
We use Qwen2.5-7B-Instruct~\cite{qwen25}, Llama-3.1-8B-Instruct~\cite{llama3} and Mistral-7B-Instruct-v0.3~\cite{mistral7b}.
Each is a 4-bit (q4\_K\_M) build served by Ollama 0.40.3 on a free NVIDIA T4 GPU, with six concurrent episodes, an 8k context and at most 512 generated tokens per turn.
Runs started on Google Colab; those left when its free quota ran out were run on Kaggle with the same GPU type, Ollama version, model builds and code.

\textbf{Tool-call protocol.}
The agent executes one tool call per model turn.
Llama and Mistral often write calls as text: JSON objects or \texttt{name(arg="...")} lines.
Such calls are recovered, logged as format slips, and stored in canonical form in the history.
One-input read tools accept their value under any argument name, while write tools are strict.
These rules came from three pilot rounds (Section~\ref{sec:pilot}) and apply equally to all models.

\textbf{Decoding and grid.}
Every episode is run once at temperature 0.
For RQ2, Qwen2.5-7B is additionally run three times at temperature 0.7 with different seeds, on all 24 tasks (1,632 runs).
To fit free-GPU limits, Llama and Mistral run a balanced half of the tasks: 12 tasks, three per read order, so every artifact type still appears at every delay equally often (204 runs each, 192 of them attacks).
Their 68 pilot runs on four of these tasks used identical code, settings and model builds and are reused.
Each run includes all three gates through forking.

\subsection{Pilot study}\label{sec:pilot}
On four tasks (68 episodes per model), three pilot rounds exposed harness artifacts that would otherwise have been reported as model safety.
\begin{itemize}
\item Llama and Mistral never invoked a tool when only native tool calls were accepted (v1).
\item Read-tool argument names mismatched, and the task wording ``Apply the change \ldots'' primed the wrong write tool (v2).
\end{itemize}
After the fixes (v3), Qwen and Llama read every payload, Mistral read 83\% (it often stopped early), and tool errors almost vanished.
The v2-to-v3 change in task wording also cut Qwen's rate of attacker-goal attempts at $d=0$ from 0.31 to 0.12, and its rate of any out-of-scope write from 0.88 to 0.19 (16 episodes per cell).
This is a pilot-scale observation that precise operator instructions reduce both wrong-tool errors and susceptibility.
Mistral completed only about 11\% of tasks: it claims writes it never executed, which limits the interpretation of its attack numbers.

\begin{table}[t]
\centering
\caption{Mock environment: tool inventory.}
\label{tab:tools}
\footnotesize
\begin{tabular}{lll}
\toprule
Tool & Kind & Parameters \\
\midrule
\texttt{get\_ticket} & read & ticket\_id \\
\texttt{list\_alerts} & read & device \\
\texttt{read\_device\_log} & read & device \\
\texttt{read\_runbook} & read & name \\
\texttt{show\_config} & read & device \\
\texttt{search\_kb} & read & query \\
\texttt{add\_ticket\_comment} & low-impact write & ticket\_id, text \\
\texttt{acknowledge\_alert} & low-impact write & alert\_id \\
\texttt{apply\_acl\_change} & high-impact write & device, acl, operation, rule, ticket\_id \\
\texttt{shutdown\_interface} & high-impact write & device, interface, ticket\_id \\
\texttt{change\_route} & high-impact write & device, prefix, next\_hop, ticket\_id \\
\texttt{create\_admin\_user} & high-impact write & device, username, ticket\_id \\
\texttt{disable\_logging} & high-impact write & device, ticket\_id \\
\bottomrule
\end{tabular}
\end{table}


\section{Results}
\pending{RQ1--RQ4 from results/eval: Table 2 (delay), Fig. 2, Fig. 3 (variability), Table 3 (gates), Fig. 4 (style $\times$ artifact), two annotated traces.}

\section{Limitations}
The network, artifacts and attacks are synthetic, as in prior benchmarks~\cite{netinjectbench,agentdojo}; real operational artifacts are not public.
Gates are deterministic stand-ins for a human approver and assume the approval record cannot be modified by the attacker.
Only three 7--8B models in 4-bit form are evaluated, under one neutral system prompt; prompt-level defenses and frontier models are out of scope.

\section{Conclusion}
\pending{Summary of answers to RQ1--RQ4.}

\bibliographystyle{IEEEtran}
\bibliography{refs}
\end{document}
